We showed that a conditional GAN can learn a practical approximation of the MFTR fading envelope distribution when conditioned on one of its parameters. Within the training range ($\mu\in[1,9]$), the generated conditional samples closely match the theoretical MFTR reference in both the QQ plots and distribution shape.

Generalization beyond the training range is more fragile: while median errors can remain comparable, the worst-case discrepancies increase for extrapolation (notably at larger $\mu$ values beyond the training upper limit). Finally, training curves (losses and discriminator probabilities) stabilize early and remain nearly constant, indicating that these diagnostics alone are not sufficient to judge the final conditional fit quality.
